\subsection[Challenge 10: Axion inflation with Chern-Simons]{Challenge 10: Axion inflation with Chern-Simons}
\textbf{Problem ID:} \texttt{Challenge\_10\_main}\\
\textbf{Primary inferred discipline:} Gravitation, Cosmology \& Astrophysics\\
\textbf{Secondary inferred disciplines:} High Energy \& Nuclear Physics\\
\textbf{Python implementation:} \texttt{challenges/10/answer.py}

\subsubsection*{Problem}
\paragraph*{Problem setup}


The gravitational action can be reformulated in the first-order form as a function of the tetrad $(e^A)$ and spin-connection variables $(\omega^{AB})$. Both of these are 1-forms on the \textcolor{blue}{$(3+1)$-dimensional} manifold $\mathcal{M}$. In this formalism, the curvature 2-form is
\begin{align}
R^{AB} =d\omega^{AB} + \omega^A{}_C\wedge\omega^{CB}\,.
\end{align}

\textcolor{blue}{Start with the Einstein-Hilbert action in first-order Palatini form,}
\begin{align}
{\color{blue}
\mathcal{S}_{EH} = -\frac{M_{Pl}^{2}}{4}\int \epsilon_{ABCD}\, e^A\wedge e^B\wedge R^{CD}\,,}
\end{align}
\textcolor{blue}{and add an action term $(\mathcal{S}_{\vartheta})$ for a single pseudoscalar, $\vartheta$, with an as-yet unspecified potential $V(\vartheta)$,}
\begin{align}
{\color{blue}
\mathcal{S}_{\vartheta} = \int \left[\, \frac{1}{2}\star d\vartheta\wedge d\vartheta \;-\; \frac{1}{4!}\,\epsilon_{ABCD}\, e^A\wedge e^B\wedge e^C\wedge e^D\; V(\vartheta) \right],}
\end{align}
\textcolor{blue}{where $\star$ denotes the Hodge dual and $\tfrac{1}{4!}\epsilon_{ABCD}e^A\wedge e^B\wedge e^C\wedge e^D = \sqrt{-g}\,d^4x$ is the spacetime volume form, so that $\mathcal{S}_{\vartheta}$ is the canonically normalised scalar action.}

We add a Chern-Simons action, which we write as
\begin{align}
\mathcal{S}_{\rm CS}=\frac{\alpha}{4}\int d\vartheta\wedge \left( \omega^{AB} \wedge d\omega_{AB} +\frac{2}{3}\omega^A{}_B\wedge\omega^B{}_C\wedge\omega^C{}_A\right),
\end{align}
where we have dropped a boundary term. 

\textcolor{blue}{The most general homogeneous and isotropic torsion 2-form contains, in addition to the two functions appearing below, an anisotropic piece $h^i{}_j(t)\,e^0\wedge e^j$ and a term $f^i{}_l\,\epsilon^l{}_{jk}\,e^j\wedge e^k$. Imposing the torsion equation of motion,}
\begin{align}
{\color{blue}
\epsilon_{ABCD}\,T^C\wedge e^D = \frac{\alpha}{M_{Pl}^{2}}\, d\vartheta\wedge R_{AB}\,,}
\end{align}
\textcolor{blue}{together with the antisymmetry of the spin connection, forces those extra components to vanish. The ansatz below is therefore \emph{not} an assumption about the most general admissible torsion: it is the on-shell form of the solution,}
\begin{align}
T^0 =  0,
\\
T^i = h(t)e^0\wedge e^i - \phi(t)\epsilon^i_{jk} e^j \wedge e^k.
\end{align}

We split the spin connection into ''Torsion free" and ''Torsion full" parts:

\begin{equation}
\omega^{IJ} = \bar{\omega}^{IJ} + \tilde{\omega}^{IJ}
\end{equation}

Assume a \textcolor{blue}{spatially flat} FRW geometry\textcolor{blue}{, with signature $(+,-,-,-)$,}
\begin{align}
{\color{blue}
g_{\mu\nu}\,dx^{\mu}dx^{\nu} = dt^{2} - a^{2}(t)\,\delta_{ij}\,dx^{i}dx^{j}\,,
\qquad\text{i.e.}\qquad
e^{0} = dt\,,\quad e^{b} = a(t)\,dx^{b}\,,}
\end{align}
\textcolor{blue}{where} the scale factor is denoted by $a(t)$, $t$ is the cosmic time and the Hubble parameter is \textcolor{blue}{$H(t)\equiv\dot a(t)/a(t)$, with $\dot{\;}\equiv d/dt$.}

\textcolor{blue}{Throughout we work in natural units, $\hbar = c = 1$, and set $M_{Pl} = 1$, where $M_{Pl} = 1/\sqrt{8\pi G}$ is the reduced Planck mass and $G$ is the gravitational constant.}

\paragraph*{Main problem}


Use these values in the equations of motions: $\alpha = 0.0001$, $V = \frac{1}{2}m\vartheta^2$, $m = 10^{-6}$, $\vartheta[t = 0] = 15$ and $\dot{\vartheta}[t = 0] = 0.1$\textcolor{blue}{.} Using these values, give me the number of e-folds achieved at $t = 25000$.
